\section{System Overview}\label{sec:overview}

Bellwether is a standalone engine for folders of Parquet files; CSV and
JSON-lines files are read by treating fixed byte-range chunks as row groups, so
nothing above the source layer sees the difference (\doc{format.md}). It
explains why a keyed query read more than it returned, proposes a change from
exact measurement, and attests that the change returns the same rows
(\doc{architecture.md}). It does not execute queries. It emits a scan plan, a
superset of the row groups that hold the queried keys together with a residual
predicate, and an engine such as DuckDB, DataFusion, or Spark reads the rows
(\adr{011}).

\subsection{The boundary and the loop}\label{sec:loop}

One sentence governs the design (\doc{architecture.md}): \emph{the agent may
observe and propose; only the engine may measure, verify, and apply; and
nothing crosses from the first side to the second without a content-addressed
record of why.} Figure~\ref{fig:loop} shows the four stages.
(1)~\emph{Observe}: an agent or a human reads Evidence, such as fanout per
column (the files and row groups one value touches, at p50, p95, and max),
lints from the rule catalog, a workload summary, or a scan plan; observe tools
are read-only, cache-backed, make no network calls, and have a 100\,ms budget.
(2)~\emph{Propose}: the agent writes a Proposal, which is a key specification,
a layout, or a rewrite described as data (before shape, after shape, invariant).
The Proposal type has no score field, so a score cannot be injected
(\doc{threat-model.md}). (3)~\emph{Score and verify}: the engine replays the
recorded workload under the current configuration and under the proposal and
writes a Replay of per-run costs (Section~\ref{sec:whatif}); it then builds the
candidate, runs the job both ways, reconciles the outputs, and writes a
Certificate (Section~\ref{sec:cert}). (4)~\emph{Apply and attest}: the engine
refuses unless the Certificate passed and was computed against the current
manifest, applies the change, and only then writes a Receipt that links the
chain and records the manifest sequence before and after (invariant~2).

The command-line interface and the Model Context Protocol (MCP)
server~\cite{mcpspec} share one code path, so humans and agents see identical
evidence (\adr{010}). MCP tools come in three tiers that follow the stages,
\code{observe.*}, \code{propose.*}, and \code{apply.*}, and \code{apply.*}
accepts only a Certificate id.

\begin{figure*}[t]
  \centering
  % The four-stage loop and the trust boundary (docs/architecture.md).
% Redrawn from the README's loop figure (now docs/assets/readme/loop-light.svg,
% drawn by scripts/render_figures.py). Used in a figure* environment.
\begin{tikzpicture}[
    font=\footnotesize,
    stage/.style={draw, thick, rounded corners=2pt, text width=3.15cm,
                  minimum height=2.05cm, align=center, inner sep=3pt},
    agent/.style={stage, fill=black!4},
    engine/.style={stage, fill=black!12},
    rec/.style={draw, rounded corners=1pt, fill=white, inner sep=2.5pt,
                font=\footnotesize\sffamily},
    tier/.style={font=\scriptsize\ttfamily, text=black!75},
    flow/.style={-{Stealth[length=2.2mm]}, thick},
    walk/.style={-{Stealth[length=1.6mm]}, semithick, black!70},
    drop/.style={-{Stealth[length=1.6mm]}, semithick, densely dotted},
  ]
  % Stages.
  \node[agent]  (obs) at (0.0,0) {\textbf{1\quad Observe}\\[2pt]
      agent or human\\ read-only, cache-backed,\\ no network, no writes};
  \node[agent]  (pro) at (3.75,0) {\textbf{2\quad Propose}\\[2pt]
      agent or human\\ a candidate change;\\ no score field exists};
  \node[engine] (sco) at (9.05,0) {\textbf{3\quad Score + verify}\\[2pt]
      engine only\\ exact replay of recorded\\ runs; reconcile outputs};
  \node[engine] (app) at (12.8,0) {\textbf{4\quad Apply + attest}\\[2pt]
      engine only\\ refuses without a passing\\ current Certificate};

  % Interface tiers.
  \node[tier, above=2pt of obs] {observe.*};
  \node[tier, above=2pt of pro] {propose.*};
  \node[tier, above=2pt of sco] {engine internal};
  \node[tier, above=2pt of app] {apply.* (Certificate id only)};

  % Flow between stages.
  \draw[flow] (obs) -- (pro);
  \draw[flow] (pro) -- (sco);
  \draw[flow] (sco) -- (app);

  % Trust boundary.
  \draw[dashed, very thick] (6.4,1.55) -- (6.4,-1.85);
  \node[font=\footnotesize\bfseries, fill=white, inner sep=1pt] at (6.4,1.75)
      {trust boundary};
  \node[font=\scriptsize, fill=white, inner sep=1pt, align=center] at (6.4,-1.45)
      {agent authority ends};

  % Ledger band and the records each stage writes.
  \node[draw, rounded corners=2pt, fill=black!2, minimum width=16.6cm,
        minimum height=0.95cm] (ledger) at (6.4,-2.45) {};
  \node[font=\scriptsize, anchor=north] at (ledger.south)
      {ledger: append-only; each record is named by the BLAKE3 hash of its canonical JSON};
  \node[rec] (ev)   at (0.0,-2.45)  {Evidence};
  \node[rec] (pr)   at (3.75,-2.45) {Proposal};
  \node[rec] (rp)   at (8.25,-2.45) {Replay};
  \node[rec] (ce)   at (9.95,-2.45) {Certificate};
  \node[rec] (rc)   at (12.8,-2.45) {Receipt};
  \draw[drop] (obs.south) -- (ev.north);
  \draw[drop] (pro.south) -- (pr.north);
  \draw[drop] (sco.south) ++(-0.8,0) -- (rp.north);
  \draw[drop] (sco.south) ++(0.9,0) -- (ce.north);
  \draw[drop] (app.south) -- (rc.north);

  % Order in which a reviewer walks back from a Receipt.
  \draw[walk] (rc) -- (ce);
  \draw[walk] (ce) -- (rp);
  \draw[walk] (rp) -- (pr);
  \draw[walk] (pr) -- (ev);
\end{tikzpicture}

  \caption{The four-stage loop and the trust boundary, after
    \doc{docs/architecture.md}. Stages 1 and 2 are open to an agent; stages 3
    and 4 run only inside the engine. Every stage writes a record to the
    ledger. Arrows inside the ledger show the order in which a reviewer walks
    back from a Receipt to the Evidence it rests on.}
  \Description{Four boxes in a row: Observe and Propose on the agent side,
    Score plus verify and Apply plus attest on the engine side, separated by a
    dashed vertical line labeled trust boundary. Each box writes a record into
    a ledger band below: Evidence, Proposal, Replay and Certificate, Receipt.}
  \label{fig:loop}
\end{figure*}

\subsection{Records, ledger, and threat model}\label{sec:ledger}

Every record is canonical JSON named by the BLAKE3 hash~\cite{blake3} of its
bytes and stored append-only (\adr{006}); field order is part of the hash, and
every record carries a version (\doc{format.md}). The current contract,
implemented as Rust types, has five records: Evidence, Proposal, Replay,
Certificate, and Receipt, and only the engine writes the last three. Evidence
and Certificates carry provenance (input hashes, the age of the oldest input, a
creation time, and the manifest version when one exists; invariant~6), so a
stale measurement cannot pass as fresh. Because
records link by hash, a reviewer or a continuous integration gate reads one
Receipt and can check every step behind it.

The threat model claims that a compromised, confused, or hallucinating agent
can waste time by proposing bad ideas and can do nothing else
(\doc{threat-model.md}). No tool accepts a file path to write, a shell command,
or a query string to execute, because rewrites are data and never code the
engine runs. The engine never consults a model for a score, a verification
outcome, or an apply decision (\adr{009}, as amended). Content addressing makes
tampering detectable but not preventable until records are signed, and an
attacker with write access to the store is out of scope.

\subsection{Storage, crates, and invariants}\label{sec:storage}

The engine's state lives beside the data: immutable segments named by their
content hash, a manifest committed by conditional PUT, the ledger, and a
rebuildable cache (\doc{format.md}, \adr{003}). A writer that loses the
conditional PUT re-reads the manifest, unions its segments with the winner's,
and retries, so readers never block (\adr{015}). The manifest carries a
sequence number that increases by one per commit, and replays and receipts
reference manifests by it (\adr{020}). Crate dependencies point one way: the
index core and the source readers know nothing above them
(Section~\ref{sec:impl}). Table~\ref{tab:invariants} lists the six invariants;
each is paired with a test that fails if it stops holding
(\doc{invariants.md}).

\begin{table}[t]
  \caption{Invariants and the crates that enforce them, from
    \doc{docs/invariants.md}.}
  \label{tab:invariants}
  \small
  \begin{tabular}{@{}r>{\raggedright\arraybackslash}p{0.62\columnwidth}l@{}}
    \toprule
    \# & Invariant & Crate \\
    \midrule
    1 & No false negatives in pruning: a row group that holds a queried key is
        in the scan plan, on every index path. & core \\
    2 & A Receipt implies a passing Certificate computed against the manifest
        it was applied to. & ledger \\
    3 & Replay is deterministic: same manifest, workload, and proposal give a
        byte-identical Replay. & replay \\
    4 & Readers see one manifest version while a build commits. & store \\
    5 & A temporal lookup never returns a superseded row. & core \\
    6 & Evidence and Certificates carry provenance: inputs, age, and creation
        time; a Certificate also names its manifest. & ledger, profile \\
    \bottomrule
  \end{tabular}
\end{table}
